\chapter{从方程到模拟，再到实验：怎样判断一项结果可信}
\label{hall:practice}
\section{建立一个最小Hall算例前，先把未知量写全}
不要直接从一张好看的重联图开始。先声明模型，例如准中性氢、$m_e=0$、标量双温、无粒子源、给定电阻率、规定的热流和碰撞交换。演化$n,\vct u_i,p_i,p_e,\B$；由$\J=\nabla\times\B/\mu_0$和$\vct u_e=\vct u_i-\J/(en)$得到电子速度。欧姆定律给$\E$，Faraday定律推进$\B$。电子压强应沿电子流对流，不能只是复制一份沿离子流的能量式。

完整初始化至少包括：密度、两温、磁场及$\nabla\cdot\B=0$；满足总压平衡的电流片；电流由两物种如何承担；扰动的磁通幅度和波数；周期/导体/开放边界；所有耗散与热源系数。周期边界意味着出流最终会回到计算域或遇到另一片，不是无限开放空间。加入有限磁岛种子可绕过线性撕裂阶段，不能再用启动时间声称自然触发机制已预测。

如果选择全Maxwell五矩/十矩，应分别推进两套$n_s,m_sn_s\vct u_s$和能量/压力，电场也有独立演化式；高频等离子体振荡和光波刚性随之出现。这个实现与准中性Hall的变量、约束和边界不相同，不能只在Hall代码里加一个电子动量项便自称完整Maxwell双流体。

\section{归一化不是隐藏单位：完整写出一套转换}
选择参考$n_0,B_0,L_0$，并定义
\begin{equation}
V_0=\frac{B_0}{\sqrt{\mu_0n_0m_i}},\quad
 t_0=\frac{L_0}{V_0},\quad J_0=\frac{B_0}{\mu_0L_0},\quad
E_0=V_0B_0,\quad p_0=\frac{B_0^2}{\mu_0}.
\end{equation}
这里$p_0$是磁压$B_0^2/(2\mu_0)$的两倍，选哪种都可以，但无量纲系数必须一致。令帽号表示除以对应参考量，欧姆定律为
\begin{equation}
\widehat{\E}+\widehat{\vct u}_i\times\widehat{\B}
=\frac{d_{i0}}{L_0}\frac{\widehat{\J}\times\widehat{\B}}{\widehat n}
-\frac{d_{i0}}{L_0}\frac{\widehat\nabla\widehat p_e}{\widehat n}
+\frac{1}{S_0}\widehat{\J},
\qquad S_0=\frac{\mu_0L_0V_0}{\eta}.
\end{equation}
压强若改用$n_0\Theta_{e0}$归一化，压力项前还需出现$\beta_{e0}/2$。不能在归一化后忘记它原来藏在压强定义里。

\worked{读取无量纲模拟输出}
取前文实验室例并选择$L_0=d_i=0.072008$ m，则$t_0=d_i/V_A=2.088\times10^{-7}$ s，也就是$\Omega_i^{-1}$。若代码输出$t=50$、电场$E_z=0.08$、位置$x=10$，其物理值为10.44微秒、1379.5 V/m、0.7201 m。若输出$p_e=0.0161$，乘$p_0=1989.4$ Pa得到32.0 Pa；这与$\beta_e=0.0322$相符。归一化温度还必须用$p/n$换回，不能把无量纲$p$直接叫eV。

\section{网格为什么被电子尺度和Hall波限制}
空间分辨率应该由所声称的最窄物理结构决定。若要分辨$d_e$内层而只用$\Delta x=0.5d_e$，只有两格，通常不足以支持其厚度与峰值的精确结论。所谓每尺度8格也只是起点；需比较更细网格并检查实际观测量。

无质量Hall的短波频率约$\omega\sim C_Hk^2$。若取最高网格波数$k_{\max}\sim\pi/\Delta x$，显式时间积分至少需$\omega_{\max}\Delta t$在其稳定域内，因此
\begin{equation}\Delta t_H\lesssim C_{\rm time}\frac{\Delta x^2}{C_H},\end{equation}
$C_{\rm time}$依空间离散及时间积分器决定，不能作为普适常数。它比通常的对流$\Delta t\propto\Delta x$更苛刻；网格缩小一半，Hall时间步常要缩小到约四分之一。电阻扩散同样给$\Delta t\propto\Delta x^2/D_\eta$。

\worked{一个时间步数量级例子}
实验室例$C_H=V_Ad_i=2.483\times10^4$ m$^2$/s。取$\Delta x=0.5$毫米，用粗略$C_{\rm time}=0.1$作\emph{教学估计}，
\begin{equation}\Delta t_H\sim0.1\frac{(5\times10^{-4})^2}{2.483\times10^4}=1.01\times10^{-12}\,\mathrm s.\end{equation}
这相当严苛，而$\Delta x<d_e$时无质量电子模型也正在丢失物理排序。隐式/半隐式推进可以缓解稳定性限制，但并不自动给出足够的时间精度或正确的电子闭合。若推进全Maxwell，光波条件约$\Delta x/c=1.67\times10^{-12}$ s，电子等离子体频率又引入额外时间尺度；选择约化模型是一项物理和数值的共同决定。

\section{用哪些检查排除“算得像”}
最小检查应同时包含：线性波频率与解析式一致；静态Harris平衡没有无故爆发；$\nabla\cdot\B$受控；总能预算包含边界通量；网格和时间步改小后，磁通率、层宽、峰值压力及积分能量趋于稳定。若使用压强下限、密度下限或限制器，它们改变的质量和能量必须记账。

建议把“数值误差”和“物理闭合差异”分开：先固定闭合做收敛，再在相同初边值下换闭合。两个模型若用不同质量比、不同种子和不同边界，差异不能唯一归因于Hall项。比较相同时间和相同重联磁通的图各有意义；前者比较时序，后者比较相近拓扑状态。论文若在不同时间取图，必须明确。

\section{GEM基准告诉了我们什么，也没有告诉什么}
Birn等的GEM挑战用规定的Harris片和有限磁岛种子比较MHD、Hall、混合与PIC。其主设置包括质量比25、背景密度为片心参考密度的0.2、片半厚0.5$d_i$、盒$25.6d_i\times12.8d_i$。含Hall响应的多类模型在该问题中的重联磁通曲线相近，传统小均匀电阻MHD则慢得多。\citep{Birn2001}

这支持Hall尺度的物种分离对该类非线性重联的重要性，但不是一条“只加Hall就解决一切”的定理。有限种子有意避开自然线性触发；各代码的电子惯性、压力及耗散实现并不完全相同；电子层和离子喷流的宽度仍可能相差约两倍。Shay等后继比较还显示密度空洞与喷流峰值存在明显模型差异。\citep{Shay2001}

因此学习顺序应是：先能重现基准的定义和输入，再检查不止一个输出，最后才扩大到导场、非对称、三维或装置几何。这里引用已有原文证据；本次讲义没有重新运行GEM的大规模模拟。

\section{实验例：MRX怎样把Hall效应变成可测量问题}
Ren等在MRX的氢/氘、无净初始导场重联中，用磁探针数组测量面外磁场。71通道高分辨率径向阵列间距1.25毫米，报告系统误差小于10 G；另有90通道二维阵列和用于重联场的29通道阵列。二维图通过径向扫描并合并多个放电构成，不是整个平面的单次同时高分辨率成像。\citep{Ren2005}

高分辨率测得的四极场幅度约50--60 G，相比约100 G的肩部重联场并不小。较粗二维图约20--40 G，原文把较小幅度与空间分辨率联系起来。四极结构横向尺度约为$d_i$，沿出流方向约为数个$d_i$。这些是空间结构、极性和尺度的共同证据。

\worked{把原文Hall电场估计重算一次}
原文采用$B_H\sim40$ G$=0.004$ T、四极场沿出流的峰间长度$L_H\sim18$ cm，用半长0.09 m估梯度：
\begin{align}
J_R&\sim\frac{0.004}{(1.25664\times10^{-6})(0.09)}
=3.54\times10^4\,\mathrm{A/m^2},\\
\frac{J_RB_Z}{en}&\sim
\frac{(3.54\times10^4)(0.020)}{(1.60218\times10^{-19})(3\times10^{19})}
=147\,\mathrm{V/m}.
\end{align}
这与原文约140 V/m的量级及典型约100 V/m的重联电场相当。但它取的是\emph{指定上游区域的指定分量}，并不证明X点处Hall项断线，也没有完整闭合所有欧姆项。原文指出部分区域的另一个磁场分量不够准确，尚不能全面评价。

其随碰撞性的扫描中，$B_H/B_{\rm rec}$约从0.6降至0.25，同时碰撞指标$\delta/\lambda_{\rm mfp}$从约0.2升至20；误差棒来自放电差异。它支持Hall结构与较弱碰撞条件的关联，但不能排除同时存在的波动或其他耗散。推断电子流还用到了电流主要由电子承担的假设，不能改写为直接测到了完整电子速度分布。原文图1--4及式(1)是建议带着本章公式阅读的定位。

\section{读论文的一个可重复流程}
按下面顺序读，会比先看“重联率接近0.1”的结论更有效：先抄出模型的独立变量和电子破冻结项；再抄归一化、质量比、导场与尺度；第三步检查边界、种子和耗散；第四步分别读取磁通率、压力、电场分量及能量预算；最后检查实验诊断和误差。未报告的网格或热源分配不要自行补上。对初学者，能明确说出“这张图不能推出什么”也是掌握模型的标志。
